% !TeX root = ../main.tex

\chapter{Discussion and Limitations}\label{ch:discussion}

% Scope rule for this chapter: interpret what Chapter 6 reported. Do not restate the
% numbers themselves, and do not introduce results that appear nowhere else.

\section{Scope of the Claims}\label{sec:scope-of-claims}

The primary scope of this thesis is the methodological evaluation of supervised
longitudinal classification approaches for \ac{MCI}-to-\ac{AD} conversion. It does not
claim to develop or validate a normative healthy-reference framework for individualized
abnormality detection.

\outlinenote{States what the evaluated protocol licenses and what it does not, including
that the time-to-event extension was scoped out.}

\section{Capacity and Inductive Bias Under Sample Constraint}\label{sec:capacity}
\outlinenote{What the parameter-per-subject accounting implies, across the independent
comparisons that point the same way.}

\section{What Self-Supervised Pretraining Contributed}\label{sec:ssl-contribution}
\outlinenote{The reconstruction objective assessed against RQ2, across both the
\ac{DELCODE}-only and the pooled protocols.}

\section{Cohort Dependence and the Limits of Transfer}\label{sec:cohort-dependence}
\outlinenote{Cohort identity is strongly decodable from the pooled embedding and in-domain
performance differs substantially between \ac{DELCODE} and \ac{ADNI}, which together
indicate that the learned representation remains strongly cohort-dependent. Discussed as
consistent with cohort-specific representations contributing to the poor transfer, without
attributing causality, and without concluding anything about which architecture transfers
where. Also discusses what the delayed reading of the escalation threshold says about the
check as a process.}

\section{Threats to Validity}\label{sec:threats}
\outlinenote{Cohort size as the dominant constraint, the conversion label as a censored
observation, informative attrition, single-site development, and the limits of a negative
result.}

\section{Clinical Relevance}\label{sec:clinical-relevance}
